% !PRINTED 181
\weberpage{181}
6. \gesperrt{A modular function which remains unchanged under the substitution $\left(\omega, -\dfrac{1}{\omega}\right)$ and which under $(\omega, \omega+1)$ takes on the factor $e^{\mp\frac{2\pi i}{3}}$ is the product of $\gamma_3(\omega)^{\pm 1}$ with a rational function $j(\omega)$.} % [sic] German source prints gamma_3 although the stated cubic-root transformation requires gamma_2; it also omits \textquotedbl{}von\textquotedbl{} (English \textquotedbl{}of\textquotedbl{}) before j(\omega).

7. \gesperrt{A modular function which changes sign under the substitutions $\left(\omega, -\dfrac{1}{\omega}\right)$ and which under $(\omega, \omega+1)$ takes on the factor $-e^{\mp\frac{2\pi i}{3}}$ is the product of $\gamma_2(\omega)\gamma_3(\omega)^{\pm 1}$ with a rational function of $j(\omega)$.}

We shall occupy ourselves more thoroughly with the modular functions in the second part.

\begin{center}
\textbf{\S\ 55. Representation of the elliptic functions by $v$ and $\varkappa^2$.}
\end{center}

If we conceive the inversion problem of the elliptic integrals, as in \S\,51, as the task of integrating the \gesperrt{elliptic differential equations}, then it demands the representation of the three functions sn\,$v$, cn\,$v$, dn\,$v$ by the two independent variables $v$, $\varkappa^2$. This task is solved by \S\,42, but with respect to the variable $\varkappa^2$ only implicitly. One can, however, also find representations in which $\varkappa^2$ occurs explicitly, namely by series which proceed in powers of $v$ and whose coefficients are \gesperrt{rational functions of $\varkappa^2$}. The coefficients of these series cannot, it is true, be represented by a general law; they can only be computed successively. These representations are owed chiefly to \gesperrt{Weierstrass}\footnote{On \gesperrt{Weierstrass}'s theory of the elliptic functions cf. especially the ``Formeln und Lehrsätze zum Gebrauch der elliptischen Funktionen'' edited by H.\,A. \gesperrt{Schwarz}. On the series developments cf. also \gesperrt{Königsberger}, ``Vorlesungen über die Theorie der elliptischen Funktionen'', 25. Vorlesung.}.

Since the $\sigma$-functions are throughout finite and continuous functions of $u$, they can be developed in unconditionally convergent power series in $u$, and if we accordingly develop the $\sigma$-functions occurring in \S\,45, (2) in this way, we obtain a solution of our task. We therefore put
%CONT>>
